% Figura 2(a) — entidade contador100 em operação livre (Questão 1)
% Descrição: ../descricoes/figura_02a_contador100.md
% Mesma caixa da Figura 1; a diferença são as constantes ligadas às entradas.
\documentclass[border=8pt]{standalone}
\usepackage{caixapreta}
\begin{document}
\begin{tikzpicture}
  \caixa{6.2}{6.6}{contador100}

  \grupoesq{0.55}{controle}
  \entrada{1.25}{reset}{1}{síncrono}
  \entrada{2.10}{clock}{1}{1 Hz}
  \entrada{2.95}{/enable}{1}{ativo em nível baixo}
  \entrada{3.80}{load}{1}{síncrono}
  \divesq{4.35}
  \grupoesq{4.65}{valor de carga}
  \entrada{5.25}{dezload}{4}{}
  \entrada{6.05}{uniload}{4}{}

  % constantes ligadas às entradas
  \constante{2.95}{'0'}
  \constante{3.80}{'0'}
  \constante{5.25}{"0000"}
  \constante{6.05}{"0000"}
  \origem{1.25}{do sistema}
  \origem{2.10}{do sistema}

  \grupodir{0.55}{saída BCD}
  \saida{1.40}{cont\_dezena}{4}{dezena}
  \saida{2.40}{cont\_unidade}{4}{unidade}

  \rodape{10.6}{%
    \texttt{/enable = '0'} \textbf{habilita} a contagem (ativo em nível baixo);
    \texttt{load = '0'} desativa a carga, então \texttt{dezload} e \texttt{uniload} não são carregados.\\[2pt]
    \textbf{Contagem:} 00 $\rightarrow$ 01 $\rightarrow$ \ldots\ $\rightarrow$ 99 $\rightarrow$ 00.
    \textbf{Arquitetura:} dois \texttt{contador10} em cascata.}
\end{tikzpicture}
\end{document}
